\documentclass[11pt]{article}
\usepackage[margin=1in]{geometry}
\usepackage{amsmath,amssymb}
\usepackage[hidelinks]{hyperref}
\usepackage[round]{natbib}

\title{The Empirical Claim: A K-Sweep Test (Code KSW)}
\date{}

\begin{document}
\maketitle
\thispagestyle{empty}

\section{The claim, in one sentence}

When a bank recapitalises, how far above the level that triggered the
raise it chooses to land is governed by the fixed cost of raising
capital, not by the bank's risk preference --- and that relationship has
a specific, testable shape.

\section{Why this claim, and not another}

The theory paper's actual headline result is that a bank's asymmetric
attitude to capital shortfalls is legible in the thresholds it commits
to, not in the volatility of its capital. That headline was the first
candidate for an empirical test, and it was set aside --- not because it
is unresolved, but because when checked, the answer disfavours a direct
test of it. Two results from the proofs document establish this:

\begin{itemize}
\item[\textbf{SCG}] proves that the \emph{definitional} deficit-to-surplus
  variance ratio --- computed directly from the model, not from data --- is
  a constant of the regulatory cap's geometry. It carries no information
  about the preference at all.
\item A separate check applies the estimator a real data panel would use
  to simulated paths from the solved model. The \emph{estimated} ratio
  does move with the preference, monotonically, but only by about
  $0.18$--$0.22$ across the full range of the preference parameter,
  against a per-path standard deviation of $0.16$--$0.20$. The signal is
  roughly the size of its own noise: usable only aggregated across many
  banks, not diagnostic from any single one.
\end{itemize}

\noindent That makes the headline claim a poor first empirical target: it
needs many banks and a lot of data before it says anything. The K-sweep
claim (\textbf{KSW}) does not have this problem. Both sides of it --- the
gap a bank leaves above its trigger, and the fixed cost of raising
capital --- are in principle observable for a single institution.

\section{The claim in detail: two falsifiable predictions}

The model has three controls acting on a bank's capital buffer: a
continuously-adjusted risk exposure, dividends paid out at an upper
barrier, and lump-sum recapitalisation at a lower trigger. Raising capital
costs $K + (1+\kappa)\xi$ for an injection of size $\xi$ --- a fixed
component $K$ and a proportional component $\kappa\xi$. Write $x_L$ for
the recapitalisation trigger and $y_{\mathrm{post}}$ for the capital
level the bank raises to. The \textbf{gap} is
$y_{\mathrm{post}} - x_L$.

\paragraph{Prediction 1 --- the gap vanishes at zero fixed cost.}
In the solved model, at $K=0$ the trigger and the target coincide exactly
(the largest numerical deviation found was $2\times 10^{-14}$, i.e.\ zero
to machine precision). Empirically: a bank facing a negligible fixed cost
of raising capital --- a private placement to an existing large
shareholder, a regulator-arranged injection with no underwriting --- should
show a post-recapitalisation capital level at or very near the level that
triggered the raise, not systematically above it.

\paragraph{Prediction 2 --- the gap scales as the cube root of the fixed
cost.} This is the sharper and more literature-grounded prediction.
Classical impulse-control theory with a fixed transaction cost predicts
\[
  y_{\mathrm{post}} - x_L \;\propto\; K^{1/3},
\]
a result derived independently of this model, not invented for it. The
solved model reproduces it: a log-log regression of the gap against $K$
gives a slope of $0.3845$ against the theoretical $1/3$, with the ratio
of gap to $K^{1/3}$ stable between $1.26$ and $1.35$ across a fourfold
range of cost values. Empirically: rank recapitalisation episodes by a
proxy for the fixed cost involved, and the gap should grow markedly more
slowly than the cost --- proportional to its cube root, not to the cost
itself.

\paragraph{What would falsify this.} A robust, non-vanishing gap at
near-zero fixed cost; a gap that scales roughly linearly with a
fixed-cost proxy rather than sub-linearly; or no detectable clustering of
post-recapitalisation capital levels at all, which would be a more basic
failure --- it would mean the "target level" the model predicts does not
describe real recapitalisation behaviour, prior to any question about how
that target responds to cost.

\section{The data problem, stated honestly}

Testing Prediction 2 needs variation in the fixed cost across
recapitalisation episodes, and the fixed cost itself is not directly
observed. It would be a mistake to assume banks face a large, easily
quantified fixed cost: the corporate-finance literature on equity issuance
costs for ordinary firms finds the opposite of what the model assumes.
\citet{altinkilic2000} find flotation spreads are up to $85\%$ variable
cost with \emph{rising} marginal cost as issue size grows --- the reverse
of a fixed-plus-linear structure. \citet{buhner2002} find fixed costs
account for at most $14$--$24\%$ of total flotation costs in the German
market. Neither result is about bank regulatory-capital raises
specifically, which may differ (preferred stock, contingent convertibles,
regulatory pre-clearance, or crisis-period government injections could
carry a genuinely different cost structure) --- but that is a hypothesis,
not yet evidence.

The practical implication: proxy for \emph{variation} in cost structure
rather than for a fixed cost's dollar value directly. Candidates are
issuer size (smaller issuers plausibly bear proportionally larger fixed
costs), public versus private/mutual charter status, and whether the
episode falls in a stress period (2008--09, 2020) versus ordinary times,
since external capital costs are well documented to rise in stress
periods regardless of the underlying fixed-cost technology.

\section{Data source and what is needed to run this}

The source is US bank Call Reports via the FDIC BankFind API --- free,
public, quarterly, covering every FDIC-insured institution. Two schedules
carry what is needed: \textbf{RC-R} for regulatory capital ratios (the
buffer), and \textbf{RI-A}, ``Changes in Bank Equity Capital,'' which
carries both dividends declared and sale of stock/capital contributions
in the same schedule --- the exact pair of events (payout,
recapitalisation) the model is about.

Concretely, in order:

\begin{enumerate}
\item Confirm the relevant RC-R and RI-A fields are reachable and named
  as expected, and that recapitalisation events (which should be rare, by
  construction of an impulse control) occur often enough in the data to
  say anything at all.
\item Pre-register, before looking at results, what counts as a
  recapitalisation event (a materiality threshold on the stock-issuance
  line) and what counts as a dividend event, plus which fixed-cost proxy
  from \S4 will be used.
\item Construct the panel: for each recapitalisation event, the capital
  ratio the quarter before (the trigger level) and the quarter after (the
  target level); compute the gap.
\item Test Prediction 1: does the gap shrink toward zero as the fixed-cost
  proxy shrinks?
\item Test Prediction 2: regress log(gap) on log(cost proxy); compare the
  slope to $1/3$, not to $1$.
\end{enumerate}

\bibliographystyle{plainnat}
\bibliography{references}

\end{document}
